\section{From Tetris Effect to Synaptic Homeostasis}\label{sec:tetris_synaptic_homeostasis}

This phenomenon is a perfect human analogue to a core concept in machine learning: \textit{overfitting} \citep{hoel2021overfitted}. In artificial intelligence, an overfit model has essentially "memorized" its training data so specifically that it fails to generalize to new situations. A classic example is Google's DeepDream network, which, after being trained on many dog photos, began "hallucinating" dog faces in clouds and random noise \citep{fastco2015deepdream}.

The Tetris Effect is a form of cognitive overfitting: the visual system has been "overtrained" on the repetitive stimulus of falling blocks, causing it to spontaneously fire those patterns \citep{tufts2021overfitted}. This isn't limited to games; people who spend a day on a boat often feel a "phantom" rocking sensation on solid ground, as their brain has over-adapted to the pattern of the waves \citep{tufts2021overfitted}.

Sleep, then, functions as the "fix" for this overfitting. This aligns with the "overfitted brain hypothesis" (discussed in Sec 3.2), which posits that sleep and dreams serve as a biological regularization process \citep{hoel2021overfitted}. The intrusive Tetris images fade after sleep, even as the \textit{skill} of playing improves. This is believed to be a two-part process:
\begin{enumerate}
    \item \textbf{NREM (Slow-Wave) Sleep:} This phase is thought to perform synaptic "downscaling" or "pruning," weakening the over-potentiated, repetitive circuits responsible for the "ghost" images.
    \item \textbf{REM Sleep:} This phase, as described by the GANs model (Sec 3.3), injects "bizarre" and diverse "data" (dreams) to ensure the brain generalizes its new knowledge rather than getting stuck on the specific, narrow pattern of the game \citep{franceschi2024gans}.
\end{enumerate}
